\chapter{Python at Scale}\label{ch:pl:python-at-scale}

A researcher's first version of a signal walks the day's events one at a time: for each record, read its side and quantity,
update an exponentially weighted average, store it. On the laptop it costs 960 nanoseconds an event, three minutes for a day of
two hundred million messages. The same computation written as one array expression costs 14 nanoseconds an event, and streamed
through the file in chunks of ten thousand events it never holds more than a third of a megabyte, whatever the length of the
history. Nothing about the mathematics changed. What changed is how much work the interpreter does per number, how many bytes
each representation of the data costs, and what is in memory at once. This chapter measures those three costs on tick-store
data and turns the lessons into the chunked-pipeline tool of the research platform (\texttt{firm.pyscale}).

\section{What the interpreter costs}

Python is the research platform's language because it is quick to write and has the libraries; the price is that every
operation written in Python is executed by an interpreter that manipulates objects.

\begin{definition}[Interpreter overhead]\label{def:pl:python-at-scale:overhead}
The \emph{interpreter overhead}\index{interpreter overhead} of an operation is the work the Python interpreter does around the
arithmetic itself: decoding the bytecode, looking up names and attributes, checking types, creating and destroying the objects
that hold intermediate values, and counting references. For a scalar operation it is tens to hundreds of times the cost of the
arithmetic.
\end{definition}

\Cref{lst:pl:python-at-scale:loops} is the first version of the signal, twice. The first loop reads the memory-mapped records of
the \omterm{def:pl:a-tick-store-on-open-formats:store}{tick store} one at a time; every access to a field of \texttt{rec[i]} creates a record object and then a scalar object. The second converts
the two columns to Python lists once and then loops over plain integers.

\omcode{code/platforms/08-python-at-scale/python/pl_pyscale.py}{53}{73}{The signal as a Python loop: over the mapped records, and
over lists made once from the columns.}\label{lst:pl:python-at-scale:loops}

\Cref{fig:pl:python-at-scale:methods} measures them: 960 nanoseconds an event over the records, 145 over lists. The difference is
the cost of building the objects that the record access creates; what remains, 150 nanoseconds for about ten bytecode operations,
is the interpreter itself. Neither is a bug; both are the price of the language.

\begin{omfigure}
\begin{tikzpicture}
\begin{axis}[width=11cm,height=5cm,xbar,bar width=9pt,xmode=log,log origin=infty,xmin=5,xmax=3000,
  xlabel={nanoseconds per event (log scale)},y dir=reverse,
  ytick={0,1,2,3},yticklabels from table={figdata/platforms/08-python-at-scale/measured_methods.csv}{method},
  tick label style={font=\small},label style={font=\small},grid=major,grid style={gray!20},point meta=explicit,
  nodes near coords,nodes near coords style={font=\footnotesize,anchor=west,/pgf/number format/fixed,/pgf/number format/precision=1},
  table/col sep=comma]
\addplot[fill=omDef!45,draw=omDef] table[x=ns_per_event,y expr=\coordindex,meta=ns_per_event]{figdata/platforms/08-python-at-scale/measured_methods.csv};
\end{axis}
\end{tikzpicture}

\omcaption{The cost per event of one signal (the exponentially weighted average of signed volume) computed four ways on
tick-store records: a Python loop over the mapped records, a Python loop over lists, one array expression, and the same
expression streamed in chunks of 65\,536 events. Measured on a laptop (Intel Core Ultra 7 155H) under WSL2, one thread, machine
otherwise idle. Data: \texttt{bench\_pyscale.py}.}\label{fig:pl:python-at-scale:methods}
\end{omfigure}

\section{Thinking in arrays}

\begin{definition}[Array programming]\label{def:pl:python-at-scale:array}
\emph{Array programming}\index{array programming} expresses a computation as operations on whole arrays --- elementwise
arithmetic, reductions, cumulative sums, filters, sorts --- so that the interpreter is entered once per operation and the loop
over elements runs in compiled code over contiguous memory.
\end{definition}

The only hard part is recursion. An exponentially weighted average $y_t = (1-\alpha)\,y_{t-1} + \alpha\,x_t$ looks inherently
sequential, but it is a linear filter: its output is the input convolved with $\alpha(1-\alpha)^k$, and a filter routine computes
it in compiled code in one call (\cref{lst:pl:python-at-scale:ewma}). On the laptop that is 14 nanoseconds an event, 70 times
faster than the loop over records and 11 times faster than the loop over lists, and the answers are identical to the last bit:
the filter performs the same multiplications and additions in the same order.

\omcode{code/firm/pyscale/firm_pyscale.py}{39}{50}{The recursive average as a linear filter; the filter's internal state for the
first element comes from the last output of the previous chunk.}\label{lst:pl:python-at-scale:ewma}

\begin{method}[From a loop to arrays]\label{met:pl:python-at-scale:arrays}
Look for the pattern in the loop: an elementwise map (arithmetic, \texttt{where}); a running total (\texttt{cumsum}); a window
(differences of cumulative sums, or a sliding view); a linear recursion (a filter); a group-by (sort, then reduce by segments);
an as-of lookup (\texttt{searchsorted}, chapter~4). What does not fit --- a book that changes with every event, a state machine
with many branches --- is a case for compiled code (chapter~9), not for a slower loop.
\end{method}

A compiler that turns the Python loop into machine code at run time is the other route, the just-in-time compilation of One Quant
Book~13 (chapter~10); the series' environment does not install one, and chapter~9 takes the explicit route of writing the kernel in
C++ or Rust.

\section{Memory: objects, arrays and copies}

The second cost is space. A Python object carries a header (a type pointer and a reference count) before its value, and a list of
them is an array of pointers to objects scattered in memory. An array of machine numbers is the numbers.

\begin{omfigure}
\begin{tikzpicture}
\begin{axis}[width=11cm,height=4.8cm,xbar,bar width=9pt,xmin=0,xmax=72,xlabel={megabytes for one million symbols},y dir=reverse,
  ytick={0,1,2,3},yticklabels from table={figdata/platforms/08-python-at-scale/memory.csv}{representation},
  tick label style={font=\small},label style={font=\small},grid=major,grid style={gray!20},
  nodes near coords,nodes near coords style={font=\footnotesize,/pgf/number format/fixed,/pgf/number format/precision=1},
  table/col sep=comma]
\addplot[fill=omMeth!45,draw=omMeth] table[x=mb,y=k]{figdata/platforms/08-python-at-scale/memory.csv};
\end{axis}
\end{tikzpicture}

\omcaption{The same column of one million symbols (fifty distinct four-character tickers) in four representations: pandas objects
(a Python string per row), a pandas categorical (integer codes and fifty strings), an Arrow string column (offsets and bytes), and
a numpy fixed-width unicode array. Data: \texttt{fig\_pyscale.py}.}\label{fig:pl:python-at-scale:memory}
\end{omfigure}

\Cref{fig:pl:python-at-scale:memory} measures one column of symbols. As Python objects it takes 61~MB, a string object per row of
which a few bytes are the characters; as a categorical, 1.0~MB, a small integer per row and fifty strings once; as Arrow strings,
8~MB, an offset and four bytes per row; as a numpy fixed-width unicode array, 16~MB, four characters of four bytes each. The
categorical is 61 times smaller than the objects: the \omterm{def:pl:columnar-formats:encodings}{dictionary encoding} of chapter~3, in memory. Every column of a research
dataframe that holds repeated strings --- symbols, venues, conditions, sides --- is worth the conversion.

\begin{remark}[Copies]\label{rem:pl:python-at-scale:copies}
\omterm{def:pl:python-at-scale:array}{Array programming} has its own trap: every intermediate expression allocates a new array. \texttt{np.where(side == B, 1.0, -1.0) *
qty} on a day of two hundred million events allocates a boolean array, two float arrays and a product, several gigabytes of
temporaries for a result of 1.6~GB. The remedies are to compute in place where the library allows it, to choose narrow types, and,
above all, not to hold the whole day at once.
\end{remark}

\section{Chunks and out-of-core loops}

\begin{definition}[Chunked processing, peak resident memory]\label{def:pl:python-at-scale:chunked}
\emph{Chunked processing}\index{chunked processing} runs a computation over a dataset one slice at a time, carrying between slices
only the state the computation needs to continue exactly (a filter's last value, a window's tail), so that the memory it uses is
set by the chunk size and not by the size of the data. The \emph{peak resident memory}\index{peak resident memory} of a process is
the largest amount of physical memory it occupied at any moment; it, not the average, decides whether a job fits a machine.
\end{definition}

The \omterm{def:pl:a-tick-store-on-open-formats:store}{tick store}'s flat files are memory-mapped (chapter~4), so a slice of the file is a view: nothing is read until it is touched,
and nothing needs to be released. \Cref{lst:pl:python-at-scale:stream} streams the signal through the file and keeps only its
running maximum; the filter's state carries across chunk boundaries, so the result equals the whole-array result.

\omcode{code/platforms/08-python-at-scale/python/pl_pyscale.py}{152}{159}{The signal streamed through the mapped file: one chunk's
temporaries and the filter's state are all that is ever in memory.}\label{lst:pl:python-at-scale:stream}

\begin{proposition}[Chunked equals whole]\label{prop:pl:python-at-scale:exact}
If a computation over a sequence can be written as $y_t = f(y_{t-1}, x_t)$, or as a function of a window of the last $w$ inputs, then
processing the sequence in chunks, with the last output (or the last $w-1$ inputs) of each chunk passed to the next, gives the same
outputs as processing it whole.
\end{proposition}

\begin{proof}
By induction over chunks: the first output of a chunk is $f$ applied to the carried state and its first input, which is what the
whole computation applies at that position; within a chunk the recursion is the same.
\end{proof}

\begin{omfigure}
\begin{tikzpicture}
\begin{axis}[width=11cm,height=5.6cm,ymode=log,xlabel={chunk size, events},ylabel={peak memory, MB (log scale)},
  tick label style={font=\small},label style={font=\small},grid=major,grid style={gray!20},xmin=-0.8,xmax=4.4,ymin=0.01,ymax=300,
  xtick={0,1,2,3,4},xticklabels={1\,000,10\,000,100\,000,1\,000\,000,whole},ytick={0.01,0.1,1,10,100},
  yticklabels={0.01,0.1,1,10,100},
  nodes near coords,point meta=explicit,nodes near coords style={font=\footnotesize,anchor=south east,
  /pgf/number format/fixed,/pgf/number format/precision=3},table/col sep=comma]
\addplot[thick,mark=*,omDef] table[x expr=\coordindex,y=peak_mb,meta=seconds]{figdata/platforms/08-python-at-scale/measured_chunks.csv};
\end{axis}
\end{tikzpicture}

\omcaption{Peak memory of the streamed signal on a history of two million events against the chunk size, each point labelled with
its time in seconds; the last point computes the whole array at once. Memory grows with the chunk; time is lowest for chunks of
a hundred thousand events and within ten per cent of it from ten thousand to the whole array, once the per-chunk overhead is
amortised. Measured on a laptop
(Intel Core Ultra 7 155H) under WSL2, one thread, machine otherwise idle; memory counted by tracemalloc, the mapped input not
included. Data: \texttt{bench\_pyscale.py}.}\label{fig:pl:python-at-scale:chunks}
\end{omfigure}

\Cref{fig:pl:python-at-scale:chunks} is the trade-off. Chunks of a thousand events pay the interpreter once per thousand and take
0.048~s for the two million events; chunks of ten thousand take 0.025~s with a peak of 0.31~MB, and chunks of a hundred thousand,
the fastest here, 0.023~s with 2.5~MB; the whole array at once takes 0.028~s with a peak of 32~MB, 16 bytes an event. At 16 bytes an event, the series' rule of at most 1.5~GB a process allows about
94 million events at once, a small fraction of a day of a large feed; streamed, the only limit is time.

\section{Processes, threads and the interpreter lock}

\begin{definition}[Global interpreter lock]\label{def:pl:python-at-scale:gil}
The \emph{global interpreter lock}\index{global interpreter lock} of the standard Python interpreter allows one thread at a time to
execute Python bytecode; threads run Python code concurrently but not in parallel, while compiled code that releases the lock (much
of numpy, input and output) runs in parallel with them.
\end{definition}

Two threads each running a pure-Python loop of five million additions take 0.20~s against 0.11~s for one: twice the work in twice
the time, no parallelism. Two threads each sorting a numpy array of eight million values take 0.094~s against 0.086~s for one: the
sort releases the lock, and the two run largely in parallel. Python code that must use several cores therefore uses several processes.

\begin{dated}{2026-09}{An interpreter without the lock}\label{dat:pl:python-at-scale:freethreading}
PEP~703 (status Final, consulted September 2026) added a build configuration of CPython that runs Python code without the \omterm{def:pl:python-at-scale:gil}{global
interpreter lock}; starting with release 3.13 CPython supports this free-threaded build, and the official macOS and Windows installers
optionally install it. The lock remains the default. The series' environment (Python 3.10) has the lock.
\end{dated}

\begin{definition}[Process-based parallelism, serialisation cost]\label{def:pl:python-at-scale:processes}
\emph{Process-based parallelism}\index{process-based parallelism} runs work in several operating-system processes, each with its own
interpreter and lock, so that Python code runs on several cores at once. Its \emph{serialisation cost}\index{serialisation cost} is
the time spent converting data to bytes and back to move them between processes (pickling, copying through a pipe), which shared
memory or files mapped by both processes avoid.
\end{definition}

Moving a 64~MB array to a second process and summing it there takes 0.14~s pickled through a pipe and 0.03~s through shared memory,
the copy into the shared block included: the serialisation dominates the pipe, and a sum of eight million numbers is negligible
beside either. The research platform's rule follows: move work to the data, not data to the workers. Give each process a slice of a
\omterm{def:pl:a-tick-store-on-open-formats:mmap}{memory-mapped file}, a list of Parquet partitions or a shared block, and have it return a small result.

\begin{omfigure}
\begin{tikzpicture}
\begin{axis}[width=11cm,height=5cm,xbar,bar width=9pt,xmin=0,xmax=0.27,xlabel={seconds},y dir=reverse,
  xtick={0,0.05,0.10,0.15,0.20,0.25},xticklabels={0,0.05,0.10,0.15,0.20,0.25},
  ytick={0,1,2,3,4,5},yticklabels from table={figdata/platforms/08-python-at-scale/measured_parallel.csv}{task},
  tick label style={font=\small},label style={font=\small},grid=major,grid style={gray!20},
  nodes near coords,nodes near coords style={font=\footnotesize,/pgf/number format/fixed,/pgf/number format/precision=3},
  table/col sep=comma]
\addplot[fill=omThm!45,draw=omThm] table[x=seconds,y expr=\coordindex]{figdata/platforms/08-python-at-scale/measured_parallel.csv};
\end{axis}
\end{tikzpicture}

\omcaption{Moving a 64~MB array to a second process (pickled through a pipe, or placed in shared memory), and running a pure-Python
loop and a numpy sort on one thread and on two (each thread doing the whole task). Measured on a laptop (Intel Core Ultra 7 155H)
under WSL2, machine otherwise idle. Data: \texttt{bench\_pyscale.py}.}\label{fig:pl:python-at-scale:parallel}
\end{omfigure}

\section{Tutorial: four ways to compute one signal}\label{tut:pl:python-at-scale:tut}

\begin{tutorial}
\textbf{Goal.} Compute one signal over tick-store records by a loop, by arrays and in chunks, show that the answers agree, and
measure time and memory. \textbf{End state:} \cref{fig:pl:python-at-scale:methods,fig:pl:python-at-scale:memory,fig:pl:python-at-scale:chunks}.
\begin{tutsteps}
\item \textbf{History}: \texttt{pl\_pyscale.history(path, n)} writes $n$ version-1 events as a tick-store flat file.
\item \textbf{Agree}: \texttt{agree()} runs the two loops (\cref{lst:pl:python-at-scale:loops}), the array expression and the
chunked version on 20\,000 events: the differences are zero.
\item \textbf{Measure}: \texttt{bench\_pyscale.py} times the four methods, the chunk sizes and the parallel tasks.
\item \textbf{Memory}: \texttt{symbol\_memory()} for the four representations of a column.
\item \textbf{Windows}: \texttt{firm\_pyscale.rolling\_sum\_kernel(w)} carries a window's tail across chunks; check it against a whole
cumulative sum.
\end{tutsteps}
\textbf{What to change next.} Stream the signal from the week's Parquet partitions of chapter~4 instead of the flat file; split the
history between two processes by halves and combine their results, and find what state crosses the boundary.
\end{tutorial}

\section{Build: the chunked pipeline}\label{bld:pl:python-at-scale:pyscale}

\begin{build}
\bfield{Purpose} The research platform's way to run any per-event computation over histories larger than memory, and to measure what
a computation costs in time and memory before it is scheduled (chapter~13).
\bfield{Interface} \texttt{chunks(array, size)}; \texttt{run(source, kernel, state) -> (outputs, state)}, a kernel being
\texttt{kernel(chunk, state) -> (output, state)}; \texttt{ewma\_kernel(alpha)}, \texttt{rolling\_sum\_kernel(window)};
\texttt{probe(fn, *args) -> (result, seconds, peak\_bytes)}.
\bfield{Rules} Chunks are views, never copies; the carried state is the minimal one that makes the chunked result equal the whole
result, and a test proves it for every kernel; memory is measured as a peak, with the method stated.
\bfield{Acceptance tests} \texttt{code/firm/pyscale/tests/}: chunks cover the array exactly; the EWMA and rolling-sum kernels give the
whole result for every chunk size, including one and the whole length; the probe reports a peak that grows with the allocation.
\bfield{Stretch} A pool of worker processes each given a range of a mapped file (at most two here, by the resource rule); a kernel
for a trade-sign filter that needs a state machine; \omterm{def:pl:python-at-scale:chunked}{peak resident memory} from the operating system alongside tracemalloc's count.
\end{build}

\begin{omsources}
\item S.~Gross, \emph{PEP 703 -- Making the Global Interpreter Lock Optional in CPython} (Final); Python documentation,
\emph{Python support for free threading}.
\item Python documentation, \texttt{multiprocessing.shared\_memory}.
\end{omsources}

\section{Exercises}

\begin{exercise}[$\star$]\label{exo:pl:python-at-scale:1}
At 960, 145 and 14 nanoseconds an event, how long does each method take for a day of two hundred million events?
\end{exercise}

\begin{exercise}[$\star$]\label{exo:pl:python-at-scale:2}
Why does the loop over records cost about six times the loop over lists, although both run the same arithmetic?
\end{exercise}

\begin{exercise}[$\star$]\label{exo:pl:python-at-scale:3}
A dataframe has a million rows and three string columns of fifty distinct values each. Estimate its string memory as objects and as
categoricals from \cref{fig:pl:python-at-scale:memory}.
\end{exercise}

\begin{exercise}[$\star\star$]\label{exo:pl:python-at-scale:4}
Write the rolling sum of the last $w$ values as an array expression, and state the state a chunked version must carry.
\end{exercise}

\begin{exercise}[$\star\star$]\label{exo:pl:python-at-scale:5}
Two threads each ran the same pure-Python loop in twice the time of one. What would two processes have done, and what would they
have cost if each needed a 64~MB input?
\end{exercise}

\begin{exercise}[$\star\star$]\label{exo:pl:python-at-scale:6}
At 16 bytes an event for the whole-array computation, how many events fit a 1.5~GB budget at once, and what fraction of a day of the
options feed's planned 311 billion messages (chapter~2) is that?
\end{exercise}

\begin{exercise}[$\star\star\star$]\label{exo:pl:python-at-scale:7}
\emph{Coding.} Check \texttt{rolling\_sum\_kernel(50)} against a whole-array cumulative-sum version on 100\,000 random values for
chunk sizes 1, 7, 50, 1\,000 and the whole length. What is the largest difference?
\end{exercise}

\begin{exercise}[$\star\star\star$]\label{exo:pl:python-at-scale:8}
\emph{Find the flaw.} ``Our feature job is slow, so we split the day into eight parts with a thread pool of eight threads, each running
the Python feature loop on its part, and concatenate the results.''
\end{exercise}

\section{Problem: Three Minutes or Three Seconds}

\begin{problem}[{Weekend problem --- a signal over a year of events}]\label{pb:pl:python-at-scale:1}
The chapter's signal (an exponentially weighted average of signed volume), its four implementations and the measurements of
\texttt{bench\_pyscale.py}.

\medskip\noindent\textbf{Part I --- The interpreter.}
\begin{enumerate}
\item What does the interpreter do per event in the loop over records that it does not do in the array expression?
\item How many nanoseconds an event does each of the four methods cost?
\item What is the speed-up of the array expression over each loop?
\item Why are the loop's and the filter's answers identical to the last bit?
\item What kinds of loop cannot be written as array operations, and where do they go?
\end{enumerate}

\medskip\noindent\textbf{Part II --- Memory.}
\begin{enumerate}[resume]
\item How many bytes a row does a column of symbols cost as objects, as a categorical, as Arrow strings and as fixed-width unicode?
\item Why is the categorical smallest?
\item What does the whole-array computation of the signal cost in memory per event, and where does it come from?
\item How many events does a 1.5~GB budget hold at once?
\item What is \omterm{def:pl:python-at-scale:chunked}{peak resident memory}, and why does it, not the average, decide whether a job fits?
\end{enumerate}

\medskip\noindent\textbf{Part III --- Chunks and cores.}
\begin{enumerate}[resume]
\item Why does the chunked result equal the whole result, and what state is carried?
\item Which chunk size is fastest here, and what is its peak memory?
\item Why are chunks of a thousand events slower?
\item What do two threads gain on a Python loop and on a numpy sort?
\item What does it cost to move 64~MB to another process by pickling, and by shared memory?
\end{enumerate}

\medskip\noindent\textbf{Part IV --- The verdict.}
\begin{enumerate}[resume]
\item State the \emph{named result}: the speed-up of the array expression over the loops, the chunk size that minimises time and its
peak memory, and the largest history that fits the budget whole.
\item How long would a year of 250 days of two hundred million events take by each method, on one core?
\item How would you use eight cores for that year without paying the \omterm{def:pl:python-at-scale:processes}{serialisation cost}?
\item Which of the chapter's lessons would a free-threaded interpreter change, and which not?
\item In one sentence: what makes Python fast enough for \omterm{def:pl:capturing-and-storing-tick-data:tick}{tick data}?
\end{enumerate}
\end{problem}

\section{Interview questions}

\begin{interviewq}[$\star$ \iqroles{researcher, developer}]\label{iq:pl:python-at-scale:1}
Why is a Python loop over a numpy array slow, and what do you do instead?
\end{interviewq}

\begin{interviewq}[$\star\star$ \iqroles{researcher, developer}]\label{iq:pl:python-at-scale:2}
How would you compute an exponential moving average over a billion values in Python?
\end{interviewq}

\begin{interviewq}[$\star\star$ \iqroles{developer}]\label{iq:pl:python-at-scale:3}
What is the \omterm{def:pl:python-at-scale:gil}{global interpreter lock}, and when does it not matter?
\end{interviewq}

\begin{interviewq}[$\star\star$ \iqroles{researcher}]\label{iq:pl:python-at-scale:4}
Your dataframe uses 40~GB of memory. Where would you look first?
\end{interviewq}

\begin{interviewq}[$\star\star\star$ \iqroles{developer, researcher}]\label{iq:pl:python-at-scale:5}
Design a job that computes features over five years of \omterm{def:pl:capturing-and-storing-tick-data:tick}{tick data} on one 32-core machine with 128~GB of memory.
\end{interviewq}
